%% kapa: exam=1 round=31 year=2020 period=2 form=A subject=accounting subject_ko=회계학 seq=28 number=68
%% source: Q-Net 감정평가사 기출문제 (cst00309) artlSeq=5207219 file=extracted/1008/5207219/A형/A형/2020년도 제31회 감정평가사 1차 시험_2교시_A형.pdf
%% conversion: kordoc 4.24.51 -> tools/gen1.py (auto); review=proofread-sonnet; flags=none
\begin{question}{68}
상품매매기업인 (주)감평은 20x1년 초 건물(취득원가 ￦10,000,000, 내용연수 10년, 잔존가치 ￦0, 정액법 상각)을 취득하면서 다음과 같은 조건의 공채를 액면금액으로 부수 취득하였다.
\begin{framedtext}
○ 액면금액: ￦2,000,000 \par
○ 발행일: 20x1년 1월 1일, 만기 3년 \par
○ 액면이자율: 연 4\%(매년 말 이자지급) \par
○ 유효이자율: 연 8\% \par
\end{framedtext}
(주)감평이 동 채권을 상각후원가 측정(AC) 금융자산으로 분류할 경우, 건물과 상각후원가 측정(AC) 금융자산 관련 거래가 20x1년 당기순이익에 미치는 영향은? (단, 건물에 대해 원가모형을 적용하고, 계산금액은 소수점 첫째자리에서 반올림하며, 단수차이로 인한 오차가 있으면 가장 근사치를 선택한다.)
\par{\small\centering\begin{tabular}{|c|c|c|c|c|}\hline
\multirow{2}{*}{기간} & \multicolumn{2}{c|}{단일금액 ￦1의 현재가치} & \multicolumn{2}{c|}{정상연금 ￦1의 현재가치} \\ \cline{2-5}
 & 4\% & 8\% & 4\% & 8\% \\ \hline
3 & 0.8890 & 0.7938 & 2.7751 & 2.5771 \\ \hline
\end{tabular}\par}
\begin{choices}
  \item ￦143,501 증가
  \item ￦856,499 감소
  \item ￦877,122 감소
  \item ￦920,000 감소
  \item ￦940,623 감소
\end{choices}
\end{question}
\answer{3}
